% Standalone appendix. For a journal manuscript, copy the section body and
% bibliography entries into the manuscript's own document/bibliography.
\documentclass[11pt]{article}
\usepackage[a4paper,margin=25mm]{geometry}
\usepackage{amsmath,amssymb}
\usepackage[T1]{fontenc}
\usepackage{lmodern}
\usepackage[hidelinks]{hyperref}
\newcommand{\E}{\mathbb{E}}
\newcommand{\Var}{\operatorname{Var}}
\newcommand{\Prob}{\mathbb{P}}
\begin{document}

\section{Moments of perturbed NB2 observations}
\label{sec:nb2-noise-moments}

Let $C$ be a latent negative-binomial count in the NB2 parameterization,
\begin{equation}
  \E(C)=\mu,\qquad v=\Var(C)=\mu+\alpha\mu^2,\qquad
  p_0=\Prob(C=0)=(1+\alpha\mu)^{-1/\alpha}.
\end{equation}
Here $\alpha>0$ is overdispersion (\texttt{dispersion} in the code),
and R's negative-binomial size is $1/\alpha$. At $\alpha=0$, the
Poisson limit gives $v=\mu$ and $p_0=e^{-\mu}$. Observation noise $\varepsilon$
is independent of $C$, and every negative noisy response is clamped to zero:
\begin{equation}
  Y=\max(0,C+\varepsilon),\qquad m=\E(Y),\qquad V=\Var(Y).
\end{equation}
For $m>0$, define the \emph{moment-matched dispersion}
\begin{equation}
  \alpha_{\mathrm{mom}}=\frac{V-m}{m^2}.
  \label{eq:moment-dispersion}
\end{equation}
This is a descriptive moment quantity: $Y$ need not be NB, and
$\alpha_{\mathrm{mom}}$ need not equal the fitted gamma-objective dispersion.
A negative value means that no NB2 distribution matches both moments.
Without noise, $(m,V,\alpha_{\mathrm{mom}})=(\mu,v,\alpha)$ for $\mu>0$.

\subsection{Positive uniform noise}
For $\varepsilon\sim U(0,1)$, clamping has no effect. Independence and the
uniform moments $\E(\varepsilon)=1/2$ and $\Var(\varepsilon)=1/12$ give
\begin{align}
  m&=\mu+\frac12, & V&=v+\frac1{12},\\
  \alpha_{\mathrm{mom}}
   &=\frac{\alpha\mu^2-5/12}{(\mu+1/2)^2}.
\end{align}
These are standard moment identities. Ouimet~\cite{ouimet2023} studies
NB counts with this uniform jitter, particularly their median; that paper
does not validate the gamma extension as a continuous likelihood.
Also, $\lfloor Y\rfloor=C$ almost surely.

\subsection{Centered uniform noise with zero clamping}
For $\varepsilon\sim U(-1/2,1/2)$, only latent zeros can be clamped.
The following special-case calculation is our own application of
conditioning, rather than a formula attributed to the cited papers.
At $C=0$, the uniform density is one, so
\begin{equation}
  \E(Y\mid C=0)=\int_0^{1/2}u\,du=\frac18,\qquad
  \E(Y^2\mid C=0)=\int_0^{1/2}u^2\,du=\frac1{24}.
\end{equation}
At an integer $k\geq1$, clamping is inactive and the conditional moments
are $k$ and $k^2+1/12$. Averaging over $C$ therefore yields
\begin{align}
  m &= \mu+\frac{p_0}{8},\\
  \E(Y^2)&=v+\mu^2+\frac1{12}-\frac{p_0}{24},\\
  V&=v+\frac1{12}-\frac{p_0}{24}
       -\frac{\mu p_0}{4}-\frac{p_0^2}{64},\\
  \alpha_{\mathrm{mom}}
    &=\frac{\alpha\mu^2+1/12-p_0/6-\mu p_0/4-p_0^2/64}
            {(\mu+p_0/8)^2}.
\end{align}
Thus centered noise acquires a positive mean shift after clamping.
The zero mass is $p_0/2$, and rounding $Y$ recovers $C$ almost surely.
Ouimet~\cite{ouimet2023} also uses centered uniform jitter, but without
the zero-clamping operation considered here.

\subsection{Gaussian noise with zero clamping}
For $\varepsilon\sim N(0,\sigma^2)$, both latent zeros and positive counts
can cross zero. Let $\phi$ and $\Phi$ denote the standard normal density
and distribution function. We use the published left-censored normal
moments of Foi et al.~\cite{foi2008}, Eqs.~(32)--(33), rescaled to
standard deviation $\sigma$. Writing their variance as a second moment gives
\begin{align}
  a_k:=\E(Y\mid C=k)
     &=k\Phi(k/\sigma)+\sigma\phi(k/\sigma),\\
  b_k:=\E(Y^2\mid C=k)
     &=(k^2+\sigma^2)\Phi(k/\sigma)
          +k\sigma\phi(k/\sigma).
\end{align}
These conditional normal formulas are published results and are not
rederived here. Our application averages them over the NB count
probabilities $p_k=\Prob(C=k)$:
\begin{align}
  m_\sigma &=\sum_{k=0}^{\infty}p_k a_k, &
  Q_\sigma&=\sum_{k=0}^{\infty}p_k b_k,\\
  V_\sigma&=Q_\sigma-m_\sigma^2, &
  \alpha_{\mathrm{mom},\sigma}
    &=\frac{Q_\sigma-m_\sigma^2-m_\sigma}{m_\sigma^2}.
  \label{eq:gaussian-mixture}
\end{align}
Also, $\Prob(Y=0)=\sum_{k=0}^{\infty}p_k\Phi(-k/\sigma)$.
Equation~\eqref{eq:gaussian-mixture} applies separately to both proposed arms:
\begin{equation}
  \sigma=1/\sqrt{12}\quad\text{and}\quad\sigma=1.
\end{equation}
The first matches the \emph{unclamped noise variance} of both uniforms;
their final response variances are generally different. The second is a
larger perturbation. The formulas condition on the exact count: no normal
approximation to the NB distribution is used. Unlike uniform jitter,
Gaussian rounding does not generally recover $C$.

The implementation adds the clipping corrections to the exact unclipped
moments $(\mu,v+\mu^2+\sigma^2)$, summing corrections through
$k=\lceil10\sigma\rceil$. It returns bounds on omitted corrections;
the theoretical expressions above remain infinite-sum identities.

\subsection{Relation to the fitted objective}
For $\alpha>0$, the existing C++ NB2 objective uses
\begin{align}
  \ell(y;\mu,\alpha)
   ={}&\log\Gamma(y+\alpha^{-1})-\log\Gamma(\alpha^{-1})
       -\log\Gamma(y+1)\nonumber\\
     &+y\log(\alpha\mu)-(y+\alpha^{-1})\log(1+\alpha\mu).
\end{align}
This is the NB log probability on integer responses. Evaluating it at real
$y\geq0$ is a working objective, not by itself a normalized continuous
density. The software retains its existing mean caps and Poisson fallback
at $\alpha\leq10^{-8}$. The experiment reports latent $\alpha$, noisy-population
$\alpha_{\mathrm{mom}}$, and fitted $\widehat\alpha$ separately. Parameter error
and a common independent integer-count test log score provide comparisons;
raw objectives evaluated on different response transformations do not.

\begin{thebibliography}{9}
\bibitem{ouimet2023}
F. Ouimet (2023).
A refined continuity correction for the negative binomial distribution and
asymptotics of the median.
\emph{Metrika}, 86, 827--849.
\href{https://doi.org/10.1007/s00184-023-00897-2}
{doi:10.1007/s00184-023-00897-2}.

\bibitem{foi2008}
A. Foi, M. Trimeche, V. Katkovnik, and K. Egiazarian (2008).
Practical Poissonian-Gaussian noise modeling and fitting for single-image
raw-data.
\emph{IEEE Transactions on Image Processing}, 17(10), 1737--1754.
\href{https://doi.org/10.1109/TIP.2008.2001399}
{doi:10.1109/TIP.2008.2001399}.
\end{thebibliography}
\end{document}
